% Осень26!A29; нумерация внутри столбца без пустых строк.
% Источник: лекция 5, §4 (строка 131)
% Источник: лекция 5, §5 (строка 190)
% Источник: лекция 5, §6 (строка 208)
% Источник: лекция 5, §7 (строка 227)
\ticket{26}{Характеризация измеримых функций с помощью ступенчатых. Следствия}
\label{ticket:26}

\subsection*{Характеризация измеримых функций последовательностями ступенчатых}

\textbf{Напоминание (определения; см. справочник из исходного конспекта п. 0.17).} \(X=\bigsqcup_{\text{кон}} e_i\) --- \textbf{разбиение} пространства на части. Функция \(f:X\to\mathbb R\) --- \textbf{ступенчатая}, если \(X=\bigsqcup_{\text{кон}} e_i\) и \(f|_{e_i}=\mathrm{const}\); такое разбиение называется \textbf{допустимым} для \(f\). Всякая ступенчатая \(f=\sum_i c_i\chi_{e_i}\); любые две ступенчатые имеют общее допустимое разбиение; \(\max\), \(\min\), сумма, произведение конечного набора ступенчатых --- снова ступенчатые.

\textbf{Соглашение лекции:} рассматриваются только \textbf{измеримые} ступенчатые (все \(e_i\in\mathcal C\)); такая ступенчатая функция автоматически измерима.

\begin{theorem}

\emph{характеризация измеримых функций последовательностью ступенчатых.} Пусть \((X,\mathcal C,\mu)\) --- произвольное пространство с мерой и \(f:X\to[0,+\infty]\) --- неотрицательная измеримая функция (бесконечные значения разрешены). Тогда существует последовательность ступенчатых функций \(f_n\) со свойствами: \[0\le f_1\le f_2\le\dots,\qquad \lim_{n\to\infty} f_n(x)=f(x)\ \ \forall x\in X.\]

То есть всякая (неотрицательная) измеримая функция --- поточечный предел \textbf{возрастающей} последовательности ступенчатых: «ступеньки приподнимаются и прижимаются к графику \(f\)». Специфика предмета: ступеньки стоят не на интервалах, а на \textbf{измеримых множествах} («это что попало»), и функция сложная, и множества сложные --- но всё равно получается.

\end{theorem}

\textbf{Эта теорема будет зашита в определение интеграла} --- потому она важна для правильного понимания конструкции.

\subsubsection*{Идея доказательства}

Режем \textbf{область значений} (ось \(Oy\)) с шагом \(1/n\): на множестве точек, где значение \(f\) попало в «этаж» \([k/n,(k+1)/n)\), ступенчатая функция объявляется равной \textbf{нижней} границе этажа \(k/n\). Сверху добавляется последний «этаж» \(\{f\ge n\}\), где ступенька равна \(n\). При росте \(n\) сетка измельчается, ступеньки приближаются к \(f\) снизу. Проблема --- монотонность (разбиения для разных \(n\) плохо согласованы) --- решается взятием максимума.

\subsubsection*{Доказательство (конструкция)}

Фиксируем \(n\). Для \(k=0,1,\dots,n^2-1\) положим \[E_{kn}:=\Bigl\{x:\ \frac{k}{n}\le f(x)<\frac{k+1}{n}\Bigr\},\] и добавим «верхний этаж» \[E_{n^2}:=\{x:\ f(x)\ge n\}.\]

\textbf{(a) Это разбиение.} Множества \(E_{kn}\) попарно не пересекаются (неравенства отсылают точки в разные диапазоны значений), и каждая точка \(x\in X\) лежит ровно в одном из них: раз \(f\ge0\), значение \(f(x)\) попадает ровно в один из промежутков \([k/n,(k+1)/n)\), \(k=0,\dots,n^2-1\), либо \(f(x)\ge n\). Итак, \[X=\bigsqcup_{k=0}^{n^2} E_{kn}.\]

\textbf{(b) Измеримость.} \(E_{kn}=X\bigl(f\ge \tfrac kn\bigr)\cap X\bigl(f<\tfrac{k+1}{n}\bigr)\) --- пересечение множеств Лебега (одного «\(\ge\)» и одного «\(<\)»), каждое измеримо (критерий измеримости), пересечение конечное --- измеримо. Аналогично \(E_{n^2}=X(f\ge n)\).

\textbf{(c) Ступенчатая функция.} Положим \[g_n(x):=\sum_{k=0}^{n^2}\frac{k}{n}\,\chi_{E_{kn}}(x),\] т.е. \(g_n=\frac kn\) на \(E_{kn}\) (на верхнем этаже \(E_{n^2}\) значение \(\frac{n^2}{n}=n\)). Это неотрицательная \textbf{измеримая} ступенчатая функция (разбиение измеримо).

\textbf{(d) Оценки и сходимость \(g_n\to f\).} По построению на каждом куске значение ступеньки --- нижняя граница соответствующего диапазона, поэтому всюду \[0\le g_n(x)\le f(x).\] Проверим поточечную сходимость (два случая).

\emph{Случай 1: \(f(x)=+\infty\).} Тогда \(f(x)\ge n\) при любом \(n\), т.е. \(x\in E_{n^2}\) и \(g_n(x)=n\to+\infty=f(x)\).

\emph{Случай 2: \(f(x)\) конечно.} Начиная с некоторого места \(n>f(x)\), и тогда \(x\) попадает в один из «нижних» этажей \(E_{kn}\), \(k\le n^2-1\); на нём \[0\le f(x)-g_n(x)=f(x)-\frac kn<\frac1n.\] Значит, \(g_n(x)\to f(x)\) и в этом случае. (Никакая монотонность \(f\) не используется; картинка с «растущим \(n\)» --- лишь иллюстрация, \(x\) неподвижен.)

\textbf{(e) Исправление монотонности.} Последовательность \(g_n\), вообще говоря, \textbf{не} возрастает: разбиения для разных \(n\) плохо согласованы, ступеньки «прыгают». Организационная мера: \[f_n:=\max\{g_1,g_2,\dots,g_n\}.\] Максимум конечного списка ступенчатых --- ступенчатая функция (общее допустимое разбиение); \(f_n\) не убывает по \(n\) (список растёт); так как каждое \(g_j\le f\), то \[g_n\le f_n\le f .\] Поскольку \(g_n(x)\to f(x)\), по теореме о зажатой последовательности \(f_n(x)\to f(x)\) при каждом \(x\). Итак, \(0\le f_1\le f_2\le\dots\), \(f_n\to f\). \(\blacksquare\)

\subsubsection*{Замечания}

\begin{itemize}
\tightlist
\item
  Теорема доказана в \textbf{произвольном} пространстве с мерой --- не обязательно \(\mathbb R^m\) с мерой Лебега; «в любом вероятностном пространстве всё будет то же самое».
\item
  Неотрицательность существенна для конструкции; от знака освободимся в следствии 1.
\end{itemize}

\subsection*{Следствие 1: приближение произвольной измеримой функции ступенчатыми}

\textbf{Следствие 1.} Пусть \(f:X\to\overline{\mathbb R}\) --- произвольная измеримая функция. Тогда существует последовательность ступенчатых \(f_n\) с \[\lim_{n\to\infty} f_n(x)=f(x)\ \forall x,\qquad |f_n(x)|\le|f(x)|\ \forall x\ \forall n.\] (Монотонности теперь, вообще говоря, нет: там, где \(f\) положительна, приближение возрастает, где отрицательна --- убывает.)

\begin{proof}

Рассмотрим срезки \(f^+\) и \(f^-\) (измеримы, неотрицательны, \(f=f^+-f^-\), \(|f|=f^+ + f^-\)). По теореме подберём ступенчатые \[0\le h_n\le f^+,\quad h_n\to f^+;\qquad 0\le g_n\le f^- ,\quad g_n\to f^- .\] Положим \(f_n:=h_n-g_n\). Разность ступенчатых --- ступенчатая. Далее, \[f_n=h_n-g_n\ \longrightarrow\ f^+-f^-=f,\] и \[|f_n|\le h_n+g_n\le f^+ + f^- =|f|. \]

\end{proof}

\textbf{Пояснение лектора, почему всё «склеивается идеально».} Там, где \(f(x)=0\): \(f^+(x)=f^-(x)=0\), значит \(h_n(x)=g_n(x)=0\). Там, где \(f>0\): \(f^-=0\), значит \(g_n\equiv0\) в этой области, и работают только ступеньки \(h_n\); где \(f<0\) --- наоборот. Иными словами, у \(h_n\) содержательные ступеньки стоят лишь на \(\{f>0\}\) (на остальном --- одна большая нулевая ступенька), у \(g_n\) --- лишь на \(\{f<0\}\); при сложении нетривиальные ступеньки объединяются, нулевые «мягко» дополняют друг друга, и конфликтов не возникает.

\subsection*{\texorpdfstring{Следствие 2: измеримость произведения, \(0\cdot\infty=0\)}{Следствие 2: измеримость произведения, 0\textbackslash cdot\textbackslash infty=0}}

\textbf{Следствие 2.} Пусть \(f,g:X\to\overline{\mathbb R}\) измеримы. Тогда \(fg\) измерима, если доопределить умножение соглашением \[0\cdot(\pm\infty):=0.\]

\textbf{О соглашении.} Значения в \(\overline{\mathbb R}\) приводят к коллизии: \(f(x)=0\), \(g(x)=+\infty\) --- как умножать? Лектор: это «настоящий, арифметический» нуль (не предельный), и «настоящий ноль бесконечность точно победит». Двусмысленность возникает лишь с «аппроксимативным» нулем (пределом), где конкурируют скорости стремления к \(0\) и к \(\infty\); здесь же нуль константный. С этим соглашением \textbf{умножение измеримых функций определено всюду} --- «умножайте всё, что угодно».

\begin{proof}

Возьмём ступенчатые \(f_n\to f\), \(g_n\to g\) со свойствами \(|f_n|\le|f|\), \(|g_n|\le|g|\) (следствие 1). Важно: ступенчатые функции по определению \textbf{конечнозначны}. Рассмотрим поточечный предел \[f_n(x)\,g_n(x)\ \longrightarrow\ f(x)\,g(x)\ ?\]

\begin{itemize}
\tightlist
\item
  Если \(f(x),g(x)\) конечны --- обычная арифметика пределов.
\item
  Если, скажем, \(f(x)=+\infty\), \(g(x)>0\) (конечно или бесконечно): \(f_n(x)\to+\infty\), \(g_n(x)\to g(x)\), и \(f_ng_n\to+\infty=f\,g\).
\item
  \textbf{Ключевая точка:} \(f(x)=0\), \(g(x)=+\infty\). Тогда \(|f_n(x)|\le|f(x)|=0\), т.е. \(f_n(x)=0\) \textbf{точно, при всех \(n\)} (не «стремится к нулю», а равна нулю!). Поэтому \(f_n(x)g_n(x)=0\cdot g_n(x)=0\) (произведение нуля с конечным числом), и предел равен \(0=f(x)\,g(x)\) по соглашению.
\end{itemize}

Итак, \(f_ng_n\to fg\) поточечно, \(f_ng_n\) --- ступенчатые (произведение ступенчатых ступенчатое), значит измеримые; по пункту 3 теоремы §2 предел \(fg\) измерим.

\end{proof}

\subsection*{Следствие 3: измеримость суммы}

\textbf{Следствие 3.} Пусть \(f,g:X\to\overline{\mathbb R}\) измеримы и \textbf{нет точки} \(x\in X\), в которой \(f(x)\) и \(g(x)\) --- бесконечности \textbf{разных знаков} (т.е. не реализуется ситуация \(+\infty+(-\infty)\)). Тогда \(f+g\) измерима.

\textbf{Оговорка («звёздочка»).} «Рекламная версия» говорит: складывайте измеримые --- всё измеримо; звёздочка: коллизию \(\infty+(-\infty)\) разрешить нельзя никаким соглашением (в отличие от \(0\cdot\infty\)), поэтому такие точки просто исключаются условием. «Мы создаём теорию и показываем, где контрольные точки; подстраивайтесь».

\begin{proof}

Аналогично: ступенчатые \(f_n\to f\), \(g_n\to g\); тогда \(f_n+g_n\to f+g\) (в \(\overline{\mathbb R}\) сумма всюду однозначно определена --- именно здесь нужна оговорка), \(f_n+g_n\) ступенчатые, измеримость предела --- по теореме §2, п. 3.

\end{proof}
